% 第1章 绪论

\chapter{绪论}\label{ch:introduction}

\section{研究背景与动机}\label{sec:background}

工业机器人在装配、抓取、人机协同等任务中的可靠性高度依赖于关节编码器的准确性。无论是底层关节伺服、笛卡尔阻抗控制、还是上层任务规划，控制律与观测模块都以编码器读数 $\vect{q}_{\text{measured}}$ 作为反馈输入，并通过正运动学与雅可比矩阵把关节空间映射到任务空间~\cite{siciliano2010robotics}。然而真实部署中编码器与真实关节角度之间普遍存在系统性偏差 $\vect{b}$，使得 $\vect{q}_{\text{measured}} = \vect{q}_{\text{true}} + \vect{b}$。

该偏差可由出厂标定残差、温度漂移、更换部件未重标、碰撞后错位、电磁干扰等多种机制引入（详细分类与典型量级见\cref{sec:bias_sources}），其中"更换部件未重标"与"碰撞后错位"两类偏差量级可达 $0.1$\,rad 以上，足以让原本可靠的策略在执行精细任务时显著退化。针对此类偏差，工业实践通常采取两类应对：\emph{停机重新标定}~\cite{kana2022fast,gong2000nongeometric}（精度高但需中断生产、影响产线节拍）与\emph{传统容错控制}（FDI/FTC，\cite{blanke2016diagnosis}；依赖明确的故障模型且以"保持稳定"而非"完成任务"为目标，在"未知量级的隐变量偏差"这类问题上往往无能为力）。本文采取与上述两类传统方案不同的路径：\textbf{在不停机的前提下，让策略在训练阶段就见过多种偏差，从而在线适应未知的校准偏差。}对该路径的若干自然替代方案——如显式辨识偏差再补偿（仍需独立外部观测且需停机重计算雅可比）、meta-RL（在隐变量恒定假设下退化为单任务 DR）——的对比与定位见\cref{ch:related_work}。这一思路要求重新审视控制问题的形式化结构（偏差作为隐变量使 MDP 退化为 POMDP）、观测设计、训练范式与人机安全保障，这些问题构成本文的研究主题。

\section{问题陈述}\label{sec:problem_statement}

引入未知偏差 $\vect{b}$ 后，原本的 MDP 在三类相互关联的层面被破坏。

\paragraph{部分可观测性}单一偏差源沿三条独立路径产生效应——初始位置偏移（PD 控制律稳态结果偏离 $\vect{q}_{\text{home}}$）、控制偏差（任务空间力矩计算基于错误的雅可比）、感知偏差（关节读数与正运动学末端位姿同时被污染）。三种效应均使智能体无法直接观测真实关节状态，控制问题从 MDP 退化为 POMDP（\cref{ch:problem_formulation} 进一步证明：基础观测下偏差不可辨识，引入独立外部位置观测后近似可辨识）。

\paragraph{部署条件约束}视觉 sim-to-real gap 使纯仿真训练加直接迁移在真机上失效；真机数据采集成本高，可用 demo 量级 $50$–$200$ 条；\textbf{本文真机条件下}（$10$\,Hz 控制频率、单 episode $\le 30$\,s、$50$ demo 量级），稀疏奖励下 SAC online 阶段 actor 解冻后出现大幅漂移（详细复盘见\cref{subsec:hilserl_failure}）。这些约束共同决定了"仿真验证加真机模仿部署"双端范式——仿真侧用 RL 验证理论可行性，真机侧用模仿学习（Imitation Learning，IL；含行为克隆 BC 与人在回路介入修正）完成部署。

\paragraph{人机协同安全}真机部署涉及人手间歇进入工作空间（HIL 训练时操作者干预、生产部署时工人接近），需要一个备份避让策略与任务策略协同，在距离阈值触发时接管控制、避让结束后交还。备份策略本身需在不引入额外不安全因素的前提下训练，本文采用纯仿真 SAC 加域随机化的方案（HIL 训练期间的安全约束亦由该备份策略提供，与上一类约束相互关联）。

\section{论文贡献}\label{sec:contributions}

本文的核心贡献围绕标题 "Hierarchical Imitation and Reinforcement Learning for Fault-Tolerant Manipulation" 展开，分三点。

\textbf{贡献 1：编码器偏差的 POMDP 形式化与可观测性分析}（\cref{ch:problem_formulation}）。本文给出一源三效应的因果链建模、POMDP 七元组形式化，并以严格命题证明：基础观测下 $(\vect{q}_{\text{true}}, \vect{b})$ 不可辨识（\cref{prop:non_identifiable}），引入独立外部位置观测后近似可辨识（\cref{prop:identifiable}）。此为\emph{信息论意义下的可辨识性}，不直接蕴含具体策略类的可学习性；\cref{subsec:identifiability_vs_solvability} 进一步澄清"可辨识性 $\neq$ 可求解性"——即便信息充分，最优策略仍必须以 $\vect{b}$ 为条件输出补偿动作，从而推出"随机偏差训练"的必要性。

\textbf{贡献 2：分层模仿与强化学习的容错操作框架}（\cref{ch:task_policy}–\cref{ch:real_chapter}）。本文设计了 task $\oplus$ backup 双策略分层架构与 sim/real 二分的双端范式，包含：

\begin{enumerate}
  \item \textbf{共享方法层}（\cref{ch:task_policy}）——观测设计、随机偏差训练、共享网络架构、稀疏奖励；
  \item \textbf{仿真侧任务策略训练}（\cref{ch:sim_chapter}）——HIL-SERL 框架验证可观测性命题（H1–H3 实验闭环）；
  \item \textbf{真机侧任务策略部署}（\cref{ch:real_chapter}）—— BC 加 HG-DAgger 模仿迭代，避开 SAC 在小数据真机的 actor 漂移问题；
  \item \textbf{备份策略仿真训练加真机部署}——纯仿真 SAC 加 DR 训练（\cref{sec:backup_sim_train}），真机阶段由 Hierarchical Supervisor（分层监督器）三态 FSM 与任务策略硬切换（\cref{subsec:supervisor}）。
\end{enumerate}

\noindent 范式选择由（观测模态 $\times$ 数据规模 $\times$ 安全约束）三轴决策原则形式化（详见\cref{sec:task_method_motivation} 与\cref{sec:backup_paradigm}）；task / backup 双策略是该决策原则在本文系统上的两个实例。

\textbf{贡献 3：真机三任务在编码器偏差注入下的端到端验证}（\cref{sec:real_experiments}）。本文实现了 B$+$D 双注入点编码器偏差注入（\cref{subsec:bd_injection}），并在 \texttt{pickup}、\texttt{wipe}、\texttt{pickandplace} 三个真机任务上设计了 $4$ 列联合成功率矩阵评估（$2\times2$ “训练偏差 $\times$ 测试偏差”析因设计：训练无偏差/测试无偏差、训练无偏差/测试有偏差、训练有偏差/测试无偏差、训练有偏差/测试有偏差），构成本论文的核心实证主线（实验设计与 partial results 见\cref{sec:real_experiments}，完整矩阵随真机数据持续填充）。同时给出 HIL-SERL 真机失败的方法论级 negative result（详见\cref{subsec:hilserl_failure}），为真机改用 BC 加 HG-DAgger 的范式选择提供实证支撑。

\section{论文组织结构}\label{sec:organization}

本文共 7 章，按以下顺序组织：

\begin{itemize}
  \item \textbf{第 1 章 绪论}：研究背景、问题陈述、贡献声明、组织结构。
  \item \cref{ch:related_work}（相关工作）：综述 FDI/FTC、POMDP 求解、模仿学习与人在回路、安全 RL 等领域文献，定位本文方法选择。
  \item \cref{ch:problem_formulation}（编码器偏差建模与 POMDP 分析）：理论核心——一源三效应的因果链、POMDP 形式化与可观测性命题。
  \item \cref{ch:task_policy}（任务策略：方法设计与共享网络架构）：仿真与真机共享的方法层。
  \item \cref{ch:sim_chapter}（仿真：策略训练与可观测性验证）：仿真侧两条主线（HIL-SERL 任务策略加 SAC$+$DR 备份策略）与 H1–H3 实验闭环。
  \item \cref{ch:real_chapter}（真机：模仿部署与端到端验证）：真机部署完整栈与三任务联合成功率矩阵。
  \item \cref{ch:conclusion}（总结与展望）：汇总贡献并给出 future work（含真机偏差鲁棒性两步消融）。
\end{itemize}
